\CNsection{1}{Why MME Exists}

The MME is the \textbf{central control-plane entity} in the LTE/\allowbreak{}EPC architecture. It exists to:

\begin{itemize}
\tightlist
\item
  \textbf{Offload the RAN from core signaling}: In 2G/\allowbreak{}3G, the RNC/\allowbreak{}BSC handled significant core-network signaling. LTE’s flat architecture moved all NAS (Non-Access Stratum) signaling termination to the MME, simplifying eNodeB (eNB) design.
\item
  \textbf{Centralize control-plane intelligence}: Authentication, mobility management, bearer setup, and paging are all orchestrated by the MME.
\item
  \textbf{Decouple control and user planes}: The MME handles no user-plane traffic — it purely manages signaling, enabling independent scaling of CP and UP.
\end{itemize}

\Needspace{17\baselineskip}

\CNsubsection{}{Key Interfaces}

{\def\LTcaptype{none} % do not increment counter
\Needspace{13\baselineskip}
\begin{longtable}[c]{@{}cccc@{}}
\toprule\noalign{}
\rowcolor{CNink}\CNth{Interface} & \CNth{Peer Node} & \CNth{Protocol} & \CNth{Purpose} \\
\CNheadrulesep
\endhead
\bottomrule\noalign{}
\endlastfoot
S1-MME & eNB & S1-AP (SCTP) & NAS transport, handover signaling \\
S6a & HSS & Diameter & Authentication vectors, subscription data \\
S11 & S-GW & GTPv2-C & Bearer/\allowbreak{}session management \\
S10 & Other MME & GTPv2-C & Inter-MME handover, context transfer \\
S3 & SGSN & GTPv2-C & 2G/\allowbreak{}3G \ensuremath{\leftrightarrow} LTE mobility \\
\end{longtable}
}

\Needspace{14\baselineskip}

\CNsection{2}{MME Functions}

\CNchained\CNsubsection{2.1}{NAS Signaling Termination}

\begin{itemize}
\tightlist
\item
  The MME is the NAS endpoint — it terminates all NAS messages from the UE (transported transparently through eNB via S1-AP).
\item
  NAS messages are integrity-protected and ciphered between UE and MME.
\end{itemize}

\CNsubsection{2.2}{Mobility Management}

\begin{itemize}
\tightlist
\item
  Tracking Area (TA) management and TA List assignment
\item
  Handover signaling (S1-based and X2 with MME involvement)
\item
  Inter-RAT mobility (to/\allowbreak{}from 2G/\allowbreak{}3G via S3/\allowbreak{}S10)
\end{itemize}

\CNsubsection{2.3}{Authentication and Security}

\begin{itemize}
\tightlist
\item
  Retrieves authentication vectors from HSS (via S6a)
\item
  Runs EPS-AKA (Authentication and Key Agreement)
\item
  Derives NAS keys (\ensuremath{K_{\mathrm{NASint}}}, \ensuremath{K_{\mathrm{NASenc}}}) and AS keys (\ensuremath{K_{\mathrm{eNB}}})
\item
  Initiates NAS Security Mode Command
\end{itemize}

\CNsubsection{2.4}{Bearer Management}

\begin{itemize}
\tightlist
\item
  Establishes default EPS bearers during attach
\item
  Coordinates dedicated bearer activation/\allowbreak{}modification/\allowbreak{}deactivation
\item
  Interacts with S-GW (S11) and indirectly P-GW
\end{itemize}

\CNsubsection{2.5}{Paging}

\begin{itemize}
\tightlist
\item
  Pages UEs in ECM-IDLE state when downlink data arrives
\item
  Sends S1-AP Paging messages to all eNBs in the UE’s registered TA list
\end{itemize}

\CNsubsection{2.6}{Idle-Mode Reachability}

\begin{itemize}
\tightlist
\item
  Maintains UE context even when UE is IDLE (ECM-IDLE)
\item
  Knows which Tracking Areas the UE may be in
\item
  Triggers paging when data or signaling arrives
\end{itemize}

\CNsubsection{2.7}{MME Selection, Pooling, and Rebalancing}

\begin{itemize}
\tightlist
\item
  Multiple MMEs serve the same area (MME Pool)
\item
  Load balancing via relative capacity (weight factor) advertised to eNBs
\item
  Rebalancing: graceful transfer of UE contexts when adding/\allowbreak{}removing MMEs
\end{itemize}

\CNsection{3}{NAS Protocol Stack}

NAS (Non-Access Stratum) operates between UE and MME, transparent to eNB.

\CNsubsection{3.1}{EMM — EPS Mobility Management}

Manages UE registration and mobility:

\begin{itemize}
\tightlist
\item
  Attach /\allowbreak{} Detach
\item
  Tracking Area Update (TAU)
\item
  Authentication (EPS-AKA)
\item
  Security Mode Control
\item
  GUTI Reallocation
\item
  Service Request
\end{itemize}

\textbf{EMM States:}

\begin{itemize}
\tightlist
\item
  EMM-DEREGISTERED: UE not registered
\item
  EMM-REGISTERED: UE attached to network
\end{itemize}

\CNsubsection{3.2}{ESM — EPS Session Management}

Manages PDN connections and EPS bearers:

\begin{itemize}
\tightlist
\item
  PDN Connectivity Request/\allowbreak{}Accept
\item
  Bearer Resource Allocation
\item
  Bearer Modify /\allowbreak{} Deactivate
\item
  PDN Disconnect
\end{itemize}

\textbf{ESM operates piggybacked on EMM during attach} (ESM message embedded in Attach Request).

\Needspace{14\baselineskip}

\CNkeepfig{m12-attach-1}{8}

\CNsection{4}{Detailed Procedures}

\CNchained\CNsubsection{4.1}{LTE Initial Attach (Full Flow)}

\CNfigure{m12-attach-1}{LTE initial attach (full flow)\CNpart{1}{2}}

\CNfigure{m12-attach-2}{LTE initial attach (full flow)\CNpart{2}{2}}

\textbf{Key Points:}

\begin{itemize}
\tightlist
\item
  Steps 7-10: EPS-AKA + NAS security activation
\item
  Steps 13-16: GTP-C session establishment across S11/\allowbreak{}S5
\item
  The Attach Request carries a piggybacked PDN Connectivity Request (ESM)
\item
  The UE gets its IP address from P-GW (step 15)
\end{itemize}

\CNkeepfig{m12-auth-smc}{3}

\CNsubsection{4.2}{Authentication and Security Mode}

\CNfigure{m12-auth-smc}{EPS-AKA authentication and NAS security mode}

\CNkeepfig{m12-key-chain}{2}

\textbf{Key Derivation Hierarchy:}

\CNfigure{m12-key-chain}{Key derivation hierarchy}

\CNsubsection{4.3}{Tracking Area Update (TAU)}

\textbf{Triggers:}

\begin{itemize}
\tightlist
\item
  \textbf{Periodic TAU}: Timer T3412 expires (default \textasciitilde{}54~min) — confirms UE is still reachable
\item
  \textbf{Mobility-triggered TAU}: UE moves to a cell whose TAC is not in its TA List
\end{itemize}

\CNfigure{m12-tau}{Tracking area update, inter-MME and intra-MME}

\Needspace{11\baselineskip}

\CNkeepfig{m12-detach-ue}{5}

\CNsubsection{4.4}{Detach Procedure}

\CNchained\CNsubsubsection{UE-Initiated Detach}

\CNfigure{m12-detach-ue}{UE-initiated detach}

\CNkeepfig{m12-detach-net}{2}

\CNsubsubsection{Network-Initiated Detach}

\CNfigure{m12-detach-net}{Network-initiated detach}

\CNkeepfig{m12-paging}{3}

\CNsubsection{4.5}{Paging (S1 Paging Procedure)}

\CNfigure{m12-paging}{S1 paging of an idle UE}

\CNkeepfig{m12-service-request}{3}

\CNsubsection{4.6}{Service Request (UE in IDLE \ensuremath{\to} CONNECTED)}

\CNfigure{m12-service-request}{Service request (ECM-IDLE \ensuremath{\to} ECM-CONNECTED)}

\Needspace{14\baselineskip}

\CNsection{5}{Bearer Management}

\CNchained\CNsubsection{5.1}{Default Bearer Setup}

\begin{itemize}
\tightlist
\item
  Created automatically during Initial Attach (part of Create Session)
\item
  One default bearer per PDN connection (always-on, non-GBR)
\item
  QCI typically 9 (best-effort Internet) or 5 (IMS signaling)
\item
  Cannot be deactivated without disconnecting the PDN
\end{itemize}

\CNkeepfig{m12-dedicated-bearer}{3}

\CNsubsection{5.2}{Dedicated Bearer Activation}

\CNfigure{m12-dedicated-bearer}{Network-initiated dedicated bearer activation}

\CNsubsection{5.3}{Dedicated Bearer Modification}

\begin{itemize}
\tightlist
\item
  Triggered by PCRF policy change (e.g., QoS upgrade during VoLTE call)
\item
  Uses Update Bearer Request/\allowbreak{}Response across S5 \ensuremath{\to} S11 \ensuremath{\to} S1-AP
\end{itemize}

\CNsubsection{5.4}{Dedicated Bearer Deactivation}

\begin{itemize}
\tightlist
\item
  Triggered by PCRF (policy removal), P-GW, or MME
\item
  Uses Delete Bearer Request/\allowbreak{}Response
\item
  Default bearer deactivation = PDN disconnect (all bearers for that PDN removed)
\end{itemize}

\Needspace{14\baselineskip}

\CNsection{6}{MME Selection}

\CNchained\CNsubsection{6.1}{DNS-Based Selection}

\begin{itemize}
\tightlist
\item
  eNB resolves MME FQDN from TAI \ensuremath{\to} DNS returns list of MME IPs
\item
  Format: \CNcode{tac-XXXX.}\allowbreak{}\CNcode{tac-YYYY.}\allowbreak{}\CNcode{mme.}\allowbreak{}\CNcode{epc.}\allowbreak{}\CNcode{mnc\textless{}MNC\textgreater{}.}\allowbreak{}\CNcode{mcc\textless{}MCC\textgreater{}.}\allowbreak{}\CNcode{3gppnetwork.}\allowbreak{}\CNcode{org}
\end{itemize}

\CNsubsection{6.2}{Load Balancing}

\begin{itemize}
\tightlist
\item
  Each MME advertises a \textbf{Relative MME Capacity} (weight, 0-255) via S1 Setup
\item
  eNB distributes new UEs proportionally across pool members
\item
  Higher capacity value = more UEs directed to that MME
\end{itemize}

\CNsubsection{6.3}{GUMMEI — Globally Unique MME Identifier}

\CNdisplay{\mathrm{GUMMEI} = \mathrm{MCC} + \mathrm{MNC} + \text{MME Group ID (MMEGI)} + \text{MME Code (MMEC)}}

\begin{itemize}
\tightlist
\item
  MMEGI identifies the MME pool
\item
  MMEC identifies a specific MME within the pool
\item
  GUTI = GUMMEI + M-TMSI (temporary UE identity)
\end{itemize}

\Needspace{14\baselineskip}

\CNsection{7}{S1-Flex and MME Pooling}

\CNchained\CNsubsection{}{Concept}

\begin{itemize}
\tightlist
\item
  \textbf{S1-Flex}: Each eNB connects to \textbf{multiple MMEs} in a pool (not just one)
\item
  \textbf{MME Pool Area}: Geographic region served by a set of MMEs
\item
  Benefits:

  \begin{itemize}
  \tightlist
  \item
    \textbf{Redundancy}: If one MME fails, others serve UEs
  \item
    \textbf{Load sharing}: New UEs distributed by weight
  \item
    \textbf{Geo-redundancy}: Pool members can be in different data centers
  \item
    \textbf{Graceful scaling}: Add/\allowbreak{}remove MMEs without service interruption
  \end{itemize}
\end{itemize}

\CNsubsection{}{MME Overload and Rebalancing}

\begin{itemize}
\tightlist
\item
  MME sends \textbf{Overload Start} to eNBs \ensuremath{\to} eNB redirects new connections to other pool members
\item
  \textbf{Rebalancing}: MME requests UEs to re-attach to redistribute load (via implicit detach with re-attach indication)
\end{itemize}

\Needspace{26\baselineskip}

\CNsection{8}{Evolution: MME \ensuremath{\to} AMF (5G)}

{\def\LTcaptype{none} % do not increment counter
\Needspace{20\baselineskip}
\begin{longtable}[c]{@{}
  >{\raggedright\arraybackslash\hspace{0pt}}m{(\linewidth - 4\tabcolsep) * \real{0.3200}}
  >{\raggedright\arraybackslash\hspace{0pt}}m{(\linewidth - 4\tabcolsep) * \real{0.3600}}
  >{\raggedright\arraybackslash\hspace{0pt}}m{(\linewidth - 4\tabcolsep) * \real{0.3200}}@{}}
\toprule\noalign{}
\rowcolor{CNink}\CNth{Aspect} & \CNth{LTE MME} & \CNth{5G AMF} \\
\CNheadrulesep
\endhead
\bottomrule\noalign{}
\endlastfoot
Architecture & Monolithic NE & Microservice (SBI-based) \\
Protocol & GTPv2-C, Diameter & HTTP/\allowbreak{}2, JSON (SBI), NGAP \\
NAS & EMM + ESM (combined) & 5G-NAS (MM only; SM separated to SMF) \\
Session Mgmt & MME handles bearer signaling & Delegated entirely to SMF \\
Slicing & N/\allowbreak{}A & Slice-aware, routes to correct SMF per slice \\
Registration & Attach/\allowbreak{}TAU & Registration (unified) \\
State & Stateful & Designed for stateless (UDSF) \\
Interface to RAN & S1-AP & NGAP (N2) \\
Interface to UDM & Diameter S6a & Nudm (SBI) \\
\end{longtable}
}

\textbf{Key Changes:}

\begin{enumerate}
\def\labelenumi{\arabic{enumi}.}
\tightlist
\item
  \textbf{Session management decoupled}: MME handled both mobility + bearers; in 5G, AMF handles only mobility/\allowbreak{}security, SMF handles sessions
\item
  \textbf{Service-based interface}: AMF exposes Namf services, consumed by other NFs via HTTP/\allowbreak{}2
\item
  \textbf{Network slicing}: AMF performs slice selection (NSSF interaction) and routes to slice-specific SMF
\item
  \textbf{Stateless design}: UE context can be stored externally (UDSF) for resilience
\item
  \textbf{Unified registration}: No separate Attach vs TAU — single Registration procedure handles both
\end{enumerate}
